% !TEX root = .\Thesis.tex

% ================================================================================================================================
% Beginning of Introduction
% ================================================================================================================================
\chapter{INTRODUCTION}\label{ch:intro}

% ////////////////////////////////////////////////////////////////////////////////////////////////////////////////////////////
% Thesis Overview
% ////////////////////////////////////////////////////////////////////////////////////////////////////////////////////////////
\section{Thesis Overview}\label{sec:thesis_overview}

Game engines can broadly be defined as frameworks that developers can use to create games. Many engines include 2D and 3D rendering, sound effects, responsive user interfaces, lighting, and cross-platform support. Some allow the developer to easily create games through an editor via dragging and dropping while others only offer a repository of code to include as a library. A number of game engines also support the creation of particle systems, or a collection of many particles that together represent an object without well-defined surfaces. Some examples of particle systems are fireworks, snow, oceans, fire, smoke, and so on. This paper will discuss the results of studying and stressing the particle systems within ChilliSource, an open-source game engine written in C++, with the goal of understanding a complex system and exploring possible optimizations that could be made to it.

In order to aide in stressing and instrumenting the engine, an automated pong game was created. This automated game could run a series of "games" where the ball bounced around the arena without any user intervention. The various kinds of games differ in the particle systems that were attached to the ball. For example, one game could have an attached particle system that only emitted 10 particles every 2 seconds, while another game could have an attached particle system that emitted 1000 particles every 5 seconds. When any given automated game reached completion, it saved information (metrics) about the games that ran. Some metrics included the number of particles that were emitted or how long the engine spent drawing particles. The metrics generated by the automated game drove all of the studies that were performed. For example, metrics relating to how long certain areas of code ran was primarily used in studies that focused on minimizing contention in a large data structure. 

All of the studies performed are organized into two categories: optimization case studies and observation studies. The optimization case studies focused on optimizing a background task that updated information about each individual particle within a particle system (e.g. current velocity, position, size, color, etc.). These studies sought to reduce contention in a large data structure either by using multiple mutexes, data structure "sharding", or hardware-based "lock free" implementations. The observation studies, on the other hand, explored other aspects of ChilliSource's particle system. The covered topics in these observation studies are the benefits of encapsulating thread synchronization within a single object, the relationship between ChilliSource's task scheduler and its particle systems, and an interesting scenario in which particles would not visually appear.
 

% ////////////////////////////////////////////////////////////////////////////////////////////////////////////////////////////
% Thesis Organization
% ////////////////////////////////////////////////////////////////////////////////////////////////////////////////////////////
\section{Thesis Organization}\label{sec:thesis_organization}

The rest of this thesis is organized as follows:

\begin{itemize}

% Background
\item \textbf{Chapter \ref{ch:background}} Review of particle systems and the ChilliSource game engine.

% Pong Game
\item \textbf{Chapter \ref{ch:csapong_game}} Overview of the automated pong game created to facilitate case studies.

% Instrumentation
\item \textbf{Chapter \ref{ch:instrumentation}} Overview of the systems used to instrument the ChilliSource engine.

% Case Studies
\item \textbf{Chapter \ref{ch:case_studies}} Discussion of the two optimization case studies performed.

% Observations
\item \textbf{Chapter \ref{ch:observations}} Discussion of the three observation studies performed.

\end{itemize}

% ================================================================================================================================
% End of Introduction
% ================================================================================================================================
